\chapter{Conclusion and Future Work}
\label{chap:conclusion}

\section{Summary of the Work}

This thesis asked which ordinal head and which loss function serve fine-grained Arabic
readability assessment best, and whether the errors they make are near misses or far misses. The
question came from a comparison that the literature leaves open: the BAREC benchmark
\cite{elmadani2025barec} compares distance-aware losses but has no ordinal head and no
metric-aligned objective, while the two shared-task systems closest to this work add a CORAL
head \cite{nazzal2025zai} and a differentiable kappa objective \cite{trigui2025stbw} but each
does so inside its own pipeline, with its own input, its own selection rule and its own
evaluation data, so neither can be read as a measurement of what its ordinal mechanism
contributes.

The answer was pursued by building one pipeline and changing one thing at a time. The corpus,
the splits, the segmented input, the AraBERTv2 encoder and the whole training protocol were
taken from the benchmark script and held fixed, down to the library defaults that the script
does not state; the prediction head and the loss function were the only components allowed to
vary. Eleven runs were carried out on this basis: nine that differ only in head and loss, one
two-phase metric-aligned run, and one ablation on the learning rate of the CORAL head. Two more
ablations re-evaluated saved models under a tuned decision threshold and under the alternative
checkpoint selection rule, and a second track trained linear heads on frozen features over five
seeds, which both satisfies the CPU-friendly requirement of the research brief and supplies the
variance estimates that a single-seed protocol cannot.

Before any of it was interpreted, the implementation was tested against the sources rather than
assumed to match them. Every loss was checked against an independent implementation to better
than $10^{-9}$, the metric code was checked against its definitions, the split sizes and the
collapse mappings were checked against the corpus itself, and each run recorded the accelerator
count, the realised batch, the observed number of optimiser steps and a comparison between the
reloaded checkpoint and the training log. That record is what allows the numbers in
Chapter~\ref{chap:results} to be quoted at two decimal places.

\section{Findings}

\subsection{The Pipeline Reproduces the Benchmark}

Seven of the eight losses that the benchmark also reports land within 1.5 QWK points of their
published development values, with a mean absolute difference of \gapQwkMean\ points, and
cross-entropy, EMD and ordinal log loss land within \gapQwkCE, \gapQwkEMD\ and \gapQwkOLL\
points. The eighth, the regression method, is \gapQwkMSE\ points below its published value under
the benchmark's own selection rule, and the per-epoch values show why: its development QWK
reaches \qwkseldevQwkMSE\ two epochs after its development loss bottoms out, and at that epoch
its test QWK is \qwkseltestQwkMSE\ against the published 84.4. The pipeline therefore
reproduces the published work, and the one visible discrepancy is a property of the selection
rule rather than of the implementation.

\subsection{Ordinal Mechanisms Help, and They Are Equivalent among Themselves}

Every ordinal-aware objective except one soft-label setting beats the ordinal-blind baseline,
by up to three QWK points on the development set and four on the test set. The leading group
consists of four objectives: the CORAL head (\devQwkCoral\ development, \testQwkCoral\ test),
the two-phase metric-aligned schedule (\devQwkSoftqwk\ and \testQwkSoftqwk), distance-weighted
ordinal log loss (\devQwkOLL\ and \testQwkOLL) and plain regression (\devQwkMSE\ and
\testQwkMSE). The differences inside that group are a few tenths of a QWK point, they do not
keep the same order across the two splits, and each method was trained once, so this thesis
declines to name a single winner. The defensible statement is that the ordinal head and the
metric-aligned objective reach the level of the best distance-aware loss under an identical
protocol, which is what the research brief asked, and that nothing in these results shows any
of them to be better than the others.

\subsection{The Metric Rewards Removing Far Misses, Not Reducing Average Error}

Separating the error distances changed the picture that a single metric gives. The methods that
lead on QWK are not the methods that predict most levels exactly, and they are not the methods
with the smallest average error either. The lowest mean distance in the study,
\devDistSoftTwo\ levels, and the highest exact accuracy, \devAccSoftTwo\ per cent, both belong
to a soft-label configuration two QWK points behind the leaders, whose own mean distances are
the worst in the table. What the leaders do is control the tail: CORAL has the fewest errors of
three levels or more on both splits (\devFarThreeCoral\ per cent on development,
\testFarThreeCoral\ on test) and the regression method the fewest of five levels or more
(\devFarFiveMSE\ and \testFarFiveMSE), while cross-entropy is worst on both counts.

This is a usable result rather than a curiosity. A reader who needs a text placed within a
narrow band should choose an ordinal head or a regression head; a reader who needs the exact
level should choose a soft-label loss and accept two QWK points less agreement. Reporting QWK
alone, as the benchmark and the shared task do, hides a difference of nearly twenty points of
exact accuracy between methods whose QWK differs by two.

\subsection{A CORAL Head Needs Its Cut-Points Initialised, Not Learned}

Under the benchmark protocol, the 18 cut-points of the CORAL head never moved further than
\cutMoveFineCoral\ from their initial values, against a fitted spread of \cutSpreadCoral\
between the highest and the lowest. The reason is arithmetic: six epochs of \stepsPerEpoch\
steps at a learning rate of $5\times10^{-5}$ decaying to zero sum to a travel budget of
\lrSumEncoder, so a cut-point starting at zero cannot reach a value near $-6$. Initialising the
cut-points from the cumulative label frequencies of the training split, which is the one
element of the method taken from neither the benchmark nor the two shared-task papers, removes
the problem instead of spending the budget on it, and it keeps the cut-points correctly ordered
throughout training.

Two further experiments support the diagnosis from both sides. Raising the head learning rate
alone to $10^{-3}$, which multiplies the budget to \lrSumHeadlr, lets the cut-points move by
\cutMoveCoralHeadlr\ and raises development QWK by \coralHeadlrDevGain\ points and test QWK by
\coralHeadlrTestGain; and on frozen features, where the head budget is \cpuLrSum, the
initialisation makes a difference of at most \cpuZeroInitDiffMean\ QWK points, that is, none.
The initialisation is therefore a remedy for a specific arithmetic situation, and a
head-specific learning rate appears to be the better remedy where it is available.

\subsection{The Checkpoint Selection Rule Has to Be Held Fixed}

Selecting on development QWK rather than development loss is worth up to
\qwkselTestGainMax\ test QWK points, and what it is worth is not uniform: it gives
\qwkselTestGainSoftThree\ points to a soft-label setting and \qwkselTestGainCE\ to
cross-entropy, but only \qwkselTestGainCoral\ to CORAL and nothing at all to ordinal log loss,
EMD or the two-phase schedule. The rule helps exactly those methods whose loss is least aligned
with the metric. Any comparison between published systems that use different rules is therefore
partly a comparison of rules, and a like-for-like comparison has to fix the rule first. A
related, smaller finding is that the tuned decision threshold which the CORAL system reports as
a source of gain contributes \zaiTestGainCoral\ test QWK points here, once the cut-points carry
the label frequencies.

\subsection{A Frozen-Feature Study Cannot Stand In for a Fine-Tuning Study}

On frozen features the seed hardly matters, with a standard deviation over five seeds of at most
\cpuStdMaxQwk\ QWK points for every method but one, and the exception, ordinal log loss at
\cpuStdOll\ points, is itself informative about that loss. But the ranking of the losses does
not survive freezing: the rank correlation with the fine-tuned ranking is
\cpuSpearmanPooler\ on pooled features and \cpuSpearmanMean\ on mean-pooled features. The cheap
protocol answers questions about variance and about CPU-friendly training; it does not answer
the question about which loss to use when the encoder can adapt.

\section{Contributions}

\begin{enumerate}
\item A controlled comparison of nine head and loss configurations for 19-level Arabic
readability on one fixed pipeline, extending the benchmark's four losses with the CORAL head,
the four soft-label settings and a metric-aligned two-phase objective, and reporting them in one
table on both official splits.
\item A verified reproduction of the benchmark's loss comparison, with the one visible
discrepancy traced to the checkpoint selection rule.
\item An initialisation of the CORAL cut-points from training label frequencies, with a
quantitative argument from the travel budget of the schedule, a direct measurement of how far
the cut-points actually move, and two ablations that test the argument from opposite directions.
\item Two metrics of error distance, the share of errors of at least three and at least five
levels, which make the near-miss and far-miss behaviour of each loss visible, and the finding
that QWK leadership is bought by removing far misses rather than by reducing average error.
\item A quantification of how much the checkpoint selection rule and the CORAL decision
threshold are worth, which is what makes published systems with different protocols comparable.
\item A frozen-backbone track that satisfies the CPU-friendly requirement of the research
brief, supplies five-seed variance estimates, and shows that the loss ranking does not transfer
between the two regimes.
\item A reproducibility record: verified losses and metrics, asserted data, recorded batch
composition and optimiser steps, pinned corpus and encoder revisions, and a checkpoint check on
every reported number.
\end{enumerate}

\section{Limitations}

Section~\ref{sec:res-limits} states the limitations in full. The three that bound the
conclusions most tightly are: one seed per method in the fine-tuning track, so that differences
of a few tenths of a QWK point are not measurements; mixed precision, where the benchmark uses
single precision; and the absence of a zero-initialised CORAL run under the fine-tuning
protocol, which would be the direct test of the cut-point argument that the measurements and
the two ablations only support indirectly.

\section{Future Work}

\begin{enumerate}
\item \textbf{Repeat the comparable group over several seeds.} This is the single most valuable
next step, and the cheapest to justify: at eleven runs of roughly half an hour each on one
mid-range accelerator, five seeds for the nine comparable methods is a day of computation and
would turn the leading group from an observation into a measurement.

\item \textbf{Run CORAL with zero-initialised cut-points under the fine-tuning protocol.} One
run would convert the cut-point argument from a supported diagnosis into a direct result.

\item \textbf{Separate the head learning rate from the encoder learning rate for every method.}
The ablation here suggests that the protocol's single learning rate, chosen for an encoder,
underserves any head with parameters that must reach specific values. Ordinal log loss and the
soft-label family may be affected in the same way as CORAL.

\item \textbf{Give every method a two-phase schedule.} The metric-aligned objective was tested
as its authors define it, which leaves open whether a second phase of Soft-QWK would also
improve CORAL or ordinal log loss. The comparison would then be between objectives under
identical schedules.

\item \textbf{Report the error distance distribution as standard.} The far-miss metrics
introduced here cost nothing to compute and expose behaviour that QWK and mean distance both
hide; a per-level error analysis, of the kind the benchmark reports for its best model, would
extend this to the question of \emph{where} on the 19-level scale each loss fails.

\item \textbf{Extend beyond one encoder and one granularity.} The comparison is confined to
AraBERTv2 at sentence level. Whether the same ordinal mechanisms behave the same way on other
Arabic encoders, and at document level, is untested here.
\end{enumerate}

\section{Closing Remarks}

Fine-grained Arabic readability assessment is a task on which two expert annotators agree
exactly six times in ten and reach a pairwise agreement of \iaaPairwiseQwk\ QWK
\cite{elmadani2025barec}. Against that background, the best test agreement obtained in this work,
\testQwkCoralHeadlr, sits above the agreement of two annotators with each other and below the
agreement of an annotator with the adjudicated label, which is where the remaining room for
improvement lies.

The practical lesson of the thesis is narrower than a leaderboard result and more useful for
building a system. Choosing an ordinal mechanism does matter: it is worth several QWK points
over an ordinal-blind baseline, and it buys those points by not making large mistakes. But
among the ordinal mechanisms that are properly implemented and properly initialised, the
evidence here does not separate them, while two things that are usually treated as
implementation details, the checkpoint selection rule and the learning rate given to the head,
are worth as much as the choice of loss. A comparison that does not hold those fixed is
measuring the protocol rather than the method.
